\newpage \pagebreak \cleardoublepage

%- Dagstuhl workshop on decision-making under uncertainty which spoke about various application areas where uncertainty might be useful
%- We will encourage paper submissions on uncertainty visualization use cases and software for domain-specific data. Accepted work will get opportunity to present their work.
%- ParaView and VisIt barely support uncertainty vis. only FunMc2 and manadatory maxima are existing uncertainty visualizations
%- bring together experts from application domains and uncertainty visualization to write a white paper on (a) critical application-specific problems that can benefit from existing uncertainty vis techniques and (b) to understand new uncertainty vis techniques that can be investigated to meet application-specific needs.
%- We propose to conduct breakout sessions that will cluster attendees similar application areas and to identify potential uncertainty-related challenges and state-of-the-art methods. In the end, we will generate summary of all breakout sessions to discuss high-priority steps. 
%- Most relevant topics of this workshop : networking for finding common interests and possibilities for cooperation, interaction with domain experts involving analysis of their domain problems, 
%- Write one paragraph each for application, techniques and software.

emphasize emerging ideas, concepts, or technologies that are currently too nascent or too interdisciplinary for a full symposium;

bring together experts on a subject to create a report, white-paper, or proposal that requires interactive work sessions, and that is of potential interest to a substantial segment of the visualization community; and/or

encourage information flow not solely from the presenters to the audience, but rather engage the participation of all attendees, for example, through collective or small group discussions.

Possible topics may include emerging or persisting problems, development of research agendas, networking for finding common interests and possibilities for cooperation, interaction with domain experts involving analysis of their domain problems, contests on solving specific visualization problems, etc.

In choosing workshop topics to propose, please consider the list of workshops that will be pre-approved and try to avoid large overlaps. The list of the pre-approved workshops will be published as soon as it becomes available.

Note that workshops remain extremely well attended, but as new associated events and publication formats are introduced to VIS, the number of submitted papers may drop from previous years. So a more interactive format (rather than one based mainly on presentations of submitted papers) is encouraged.

Workshop proposals should include:

\begin{itemize}
\item a title

\item the contact details of the organizers

\item a brief description of the organizers’ background, related publications, and research,
the intended format (i.e., in person, remote, or hybrid),
the planned activities (e.g., mainly talks, mainly interactive sessions, a mix), including an outline schedule for the program—workshops are strongly encouraged, though not required, to consider ways to interactively engage all participants

\item the intended result and impact of the workshop,
NEW: a brief justification for why this workshop is necessary (e.g., why would the topic benefit from a more informal setting, in what ways is the topic emergent or a relatively poor fit for the main conference program, etc.)

\item a chosen measure for a successful CFP (e.g., number of submissions), possibly including an outline for a “backup policy” (e.g., changing the workshop format, or voluntarily withdrawing the workshop before acceptance notification is sent, but should not require the workshop organizers to submit emergency papers to their own workshop)

\item a statement of the organization and the development of the list of participants (intended size, detailed selection procedures, and timeline for finalizing workshop presenters),
a list of any special technology needed

\item the plan for publication (see the classes listed below),
the number of poster slots requested, if the workshop intends to feature posters 

\item the proposed dates for the call for participation, author notification, and camera-ready deadlines (author notification must be prior to the early registration deadline for IEEE VIS; for inclusion of materials on IEEE Xplore and/or the downloadable proceedings, the camera-ready deadline must be prior to August 19, 2024).

\item Because the number of workshop time slots is limited, half-day workshop proposals are strongly encouraged. However, well-justified full-day proposals will also be considered. Full-day proposals may optionally state what changes would be implemented to allow the workshop schedule to fit within a half day.

\end{itemize}
